\documentclass[11pt]{article}
\pagestyle{empty}

%\setlength{\topmargin}{.25in}

\setlength{\textwidth}{6in}
\setlength{\oddsidemargin}{.25in}
\setlength{\topmargin}{0in}
\setlength{\textheight}{8.25in}


\renewcommand{\baselinestretch}{1.1}
\usepackage{amsmath}
\usepackage{epsfig}
\usepackage{amssymb}
\usepackage{amscd}
\usepackage{times}
\usepackage{graphics}
%\usepackage{chicago}
%\usepackage{fullpage}
%\usepackage{apalike}
\newtheorem{proposition}{Proposition}
\newtheorem{corollary}{Corollary}
\newtheorem{theorem}{Theorem}
\newtheorem{lemma}{Lemma}
\newtheorem{example}{Example}
\newtheorem{claim}{Claim}
\newtheorem{assumption}{Assumption}
\newtheorem{definition}{Definition}
\def\argmax{\mathop{\rm argmax}\limits}
\def\max{\mathop{\rm max}\limits}


\begin{document}

\noindent \large Midterm 2019, ``The Mathematical Foundations of Political Methodology"

\normalsize
\ \\You may use a calculator and refer to your formula sheet. Please do not use the internet during the exam.  Show your work.

\begin{enumerate}

\item Consider the following sets of real numbers:
\begin{eqnarray}
A&=& \{x\ : 1\leq x\leq 5\} \nonumber\\
B&=& \{x\ : 0\leq x\leq 2\}\nonumber\\
C&=& \{x\ : x<0\}\nonumber
\end{eqnarray}

Describe each of the following sets of real numbers in a single sentence.
\begin{enumerate}
\item $A^c$
\bigskip\bigskip
%The numbers greater than 5 and less than 1.

\item $A\cup B$
%The numbers between 0 and 5 inclusive.
\bigskip\bigskip
\item $B\cap C^c$
% B
\bigskip\bigskip
\item $(A\cup B)\cap C$
% none.
\bigskip\bigskip

\item $(A\cup B)^c\cap C$
%C

\end{enumerate}

\vspace{.2in}

\item Find the limit of each of the following sequences as $n\rightarrow \infty$, if the limit exists.
\begin{enumerate}
\item ${{5n}\over{(n^2-5)}}$
\bigskip
\bigskip
% The limit is 0
\item ${{4n^{\frac{1}{2}}}\over{(n^2-5)}}+2$
\bigskip
\bigskip
% It's 2.
\item  $\frac{4n}{(n^{\frac{1}{2}}-4)}+2$

% It doesn't exist.
\end{enumerate}

\pagebreak
\vspace{.2in}
\item Use the $(\delta, \epsilon)$ definition of a limit to prove that $\lim_{x\rightarrow 3} 4x-11=1$
%

%\ \\Proof:  

%For any $\epsilon>0$, we can find a $\delta>0$ such that $|4x-11-1|<\epsilon$ whenever $|x-3|<\delta$.  

%Rearranging, we get $4|x-3|<\epsilon$.  $| x-3|<\frac{\epsilon}{4}$ 

% Setting $\delta=\frac{\epsilon}{4}$ proves the result. $\Box$

%\item Use the definitions of continuity and differentiability to evaluate whether or not the following function is 1) continuous and 2) differentiable at $x=0$.
%$$f(x)=x^{\frac{1}{3}}$$
\bigskip
\bigskip
\bigskip
\bigskip
\bigskip
\bigskip
\bigskip
\bigskip
\bigskip
\bigskip

\vspace{.2in}
\item Let $$f(x)=\frac{x}{1+e^x}$$
\begin{enumerate}
\item Is f(x) concave, convex or both?
% Take derivative 
%$$f'(x)=\frac{1-e^x(x-1)}{1+e^x}^2$$
%$$f''(x)=\frac{2xe^{2x}}{{(1+e^x)}^3}-\frac{xe^{x}}{{(1+e^x)}^2}-\frac{2e^{x}}{{(1+e^x)}^2}$$

%Plug in any small number, like 0.

%$$f''(0)=-\frac{2e^{0}}{{1+e^0}^2}=-\frac{2}{2}=-1<0$$

%However, plug in 5, you get a positive number

%$$\frac{10e^{10}}{{(1+e^5)}^3}-\frac{5e^{5}}{{(1+e^5)}^2}-\frac{2e^{5}}{{(1+e^5)}^2}=.02$$
%This function is positive in some bits and negative in others so it is both.
\bigskip
\bigskip
\bigskip
\bigskip
\bigskip
\bigskip
\bigskip
\bigskip
\bigskip
\bigskip

\item Find any asymptotes that $f$ may have.
% trace the plot.
%The function approaches y=x on the left and an asymptote at the right at the x-axis.
\bigskip
\bigskip
\bigskip
\bigskip
\bigskip
\bigskip
\bigskip
\bigskip
\bigskip
\bigskip
\end{enumerate}

\item Perform the following matrix operations:

\begin{enumerate}
\item Find the inverse of $\begin{bmatrix}\frac{1}{2} & \frac{1}{2}\\ 0& \frac{1}{2}
\end{bmatrix}$

%$\begin{bmatrix} 2 & -2\\ 0& 2
%\end{bmatrix}$

\pagebreak

\item Find the rank of 
$A=\begin{bmatrix} 
3& 1& 1&2\\
2& 1& -1&2\\
0& 3& 2&0\\
1& 0& 3&0
\end{bmatrix}$
\bigskip
\bigskip
\bigskip
\bigskip

%It is 4 

\item What $x$ would make the following matrix singular?
$$\begin{bmatrix}
1&7\\ x&6
\end{bmatrix}$$
%x=6/7
\end{enumerate}
\bigskip
\bigskip
\bigskip
\bigskip
\vspace{.2in}

\item If A and B are matricies, which of the following imply the others: 
\begin{enumerate}
\item AB=BA
\item Either A or B is the identity matrix.
\item B is the inverse of A.
\item $A=\frac{1}{det(B)}adj(B)$
\item A and B are symmetric.
\item A and B are square matrixes
\end{enumerate}

% a implies f
% b implies a, 
% c implies a, d and f
% d implies a and c and f
% e implies f
% f implies nothing.
\pagebreak

\item What constant $k$ solves the following equation?
$$\int_0^{10} ky^2dy=1$$
% k = \frac{3}{1000}
\bigskip
\bigskip
\bigskip
\bigskip\bigskip
\bigskip
\bigskip
\bigskip\bigskip
\bigskip
\bigskip
\bigskip\bigskip
\bigskip
\bigskip
\bigskip
\item What values of the constant $a$ satisfies the following equation?
$$\int_0^a \frac{2x}{2+x^2}dx=1$$
% a^{2} = 2 - exp(1 + log(2)), which implies a = \sqrt{2 - exp(1 + log(2))} or - \sqrt{2 - exp(1 + log(2))}
% U = 1+x^2
% du= 2x

%log(2+x^2)|_0^a  
% log(2+a^2)-log(1)=1
% 2+a^2=exp(1)
% a=\sqrt(exp(1)-2)

\end{enumerate}

\end{document}
